%------------------------------
% CLEVEREF - must be loaded after hyperref (loaded via pdfx in meta.tex)
%------------------------------
\usepackage[capitalize,noabbrev]{cleveref}
%------------------------------
% MICROTYPE
%------------------------------
\microtypecontext{spacing=nonfrench}
%------------------------------
% DIRECTORIES
%------------------------------
\setcounter{secnumdepth}{2}	% number chapter.section.subsection only (no numbered \paragraph)
\setcounter{tocdepth}{2}	% table of contents down to subsection; deeper headings are run-in, unnumbered
%------------------------------
% DATE
%------------------------------
\makeatletter
\newcommand{\monthyeardate}{\DTMenglishmonthname{\@dtm@month} \@dtm@year}
\makeatother
%------------------------------
% FIGURES, TABLES, LISTS, NUMERATIONS, COUNTER
%------------------------------
\graphicspath{{figures/}} 		% directory in which LaTeX searches for figures
\newcounter{includepdfpage}		% counter for pdfpages page references
\newlist{myList}{itemize}{1}
\setlist[myList]{leftmargin=0.0cm,nosep,label={}}
\renewcommand{\topfraction}{0.98}
\renewcommand{\bottomfraction}{0.98}
%------------------------------
% DEDICATION
%------------------------------
\newenvironment{dedication_env}	% defines the dedication environment
{\Large\begin{quote}\begin{center}\begin{em}}
{\par\end{em}\end{center}\end{quote}}
%------------------------------
% HEADER
%------------------------------
\lehead{\scshape\leftmark}		% left-even-head
\cehead{}						% center-even-head
\rehead{}						% right-even-head
\lohead{}						% left-odd-head
\cohead{}						% center-odd-head
\rohead{\scshape\rightmark}		% right-odd-head
%------------------------------
% FOOTER
%------------------------------
\lefoot[\pagemark]{\pagemark}	% left-even-foot
\cefoot{}						% center-even-foot
\refoot{}						% right-even-foot
\lofoot{} 						% left-odd-foot
\cofoot{}						% center-odd-foot
\rofoot[\pagemark]{\pagemark}	% right-odd-foot
\renewcommand*{\chapterpagestyle}{plain}
%------------------------------
% LISTS OF ...
%------------------------------
\makeatletter					%
\renewcommand*{\@pnumwidth}{2.5em}% column width for page numbers
\makeatother					%
%------------------------------
% LIST OF ABBREVIATIONS
%------------------------------
\newlength{\myabbrevlength}
\setlength{\myabbrevlength}{2cm}
\newlength{\mydescriptionwidth}
\setlength{\mydescriptionwidth}{\linewidth-2\tabcolsep-\myabbrevlength}
\newcommand{\abbrevname}{List of Abbreviations}
\NewAcroTemplate[list]{dbtalongtable}{
	\AcroNeedPackage{array}%
	\def\arraystretch{1.2}
    \acronymsmapF{%
		\AcroAddRow{%
			\leavevmode\acrowrite{short} & \acrowrite{list} \acropagefill \tabularnewline}
    }
	{\AcroRerun}%
	\acroheading%
	\acropreamble%
    \setlength\LTleft{0.5ex}    % Adjustment of left margin
    \begin{longtable}{>{\bfseries\sffamily}lp{\mydescriptionwidth}}
        \AcronymTable
    \end{longtable}
}
\acsetup{
	format/first-long=\itshape, % first appearance is set in italics
	index/use=true,				% add abbrevs to index
	make-links=true				% make links for abbrevs
}
%------------------------------
% TODOS
%------------------------------
\makeatletter
\renewcommand{\@todonotes@todolistname}{List of ToDos}
\makeatother
%------------------------------
% THEOREM STYLE
%------------------------------
\newtheoremstyle{mytheoremstyle}
{4pt}{4pt}{\itshape}{0pt}{\bfseries}{.}{5pt plus 1pt minus 1pt}
{\thmname{#1}\thmnumber{ #2}\thmnote{ (#3)}}
\newtheoremstyle{mydefinitionstyle}
{4pt}{4pt}{\normalfont}{0pt}{\bfseries}{.}{5pt plus 1pt minus 1pt}
{\thmname{#1}\thmnumber{ #2}\thmnote{ (#3)}}
\newtheoremstyle{myremarkstyle}
{4pt}{4pt}{\normalfont}{0pt}{\itshape}{.}{5pt plus 1pt minus 1pt}
{\thmname{#1}\thmnumber{ #2}\thmnote{ (#3)}}
%------------------------------
% THEOREM-LIKE ENVIRONMENTS (plain, no boxes)
%------------------------------
\theoremstyle{mytheoremstyle}
\newtheorem{theoremenv}{Theorem}[chapter]
\newtheorem{lemmaenv}[theoremenv]{Lemma}
\newtheorem{proposition}[theoremenv]{Proposition}
\newtheorem{corollary}[theoremenv]{Corollary}
\newtheorem{claim}{Claim}[chapter]
\theoremstyle{mydefinitionstyle}
\newtheorem{definition}{Definition}[chapter]
\theoremstyle{myremarkstyle}
\newtheorem{remarkenv}{Remark}[chapter]
\newtheorem{remark}[remarkenv]{Remark}
\newtheorem{observation}{Observation}[chapter]
%------------------------------
% MATH AND OPERATORS
%------------------------------
\DeclareMathOperator{\rot}{rot}	% declares the rot math operator
\newcommand*\diff{\mathop{}\!\textrm{d}}	% differential d, e.g. in integrals
%------------------------------
% PAPER-SPECIFIC COMMANDS
%------------------------------
\newcommand{\NN}{\mathbb{N}}
\newcommand{\QQ}{\mathbb{Q}}
\newcommand{\ZZ}{\mathbb{Z}}
\newcommand{\CC}{\mathbb{C}}
\newcommand{\RR}{\mathbb{R}}
\newcommand{\relu}{\text{ReLU}}

% Math operators and commands
\DeclareMathOperator{\argmax}{argmax}
\DeclareMathOperator{\argmin}{argmin}
\DeclareMathOperator{\diag}{diag}

% Custom math symbols and notation
\newcommand{\maxout}{\text{Maxout}}
\newcommand{\1}{\mathbf{1}} % indicator function or vector of ones
\renewcommand{\P}{\mathbb{P}} % probability symbol (redefines paragraph symbol)
\newcommand{\E}{\mathbb{E}} % expectation symbol
\newcommand{\indicator}[1]{\mathds{1}_{#1}} % for indicator function using dsfont package
\newcommand{\abs}[1]{\lvert#1\rvert}
\newcommand{\norm}[1]{\lVert#1\rVert}

% Additional game-theoretic notation
% \Gamma serves as the terminal state symbol (already defined in LaTeX, no macro needed)
\newcommand{\textbfplus}{\textsuperscript{+}}
\newcommand{\textbfminus}{\textsuperscript{-}}

%------------------------------
% SEMANTIC MACROS: SPLIT STREAMS & DECOMPOSITIONS
% Change notation globally here; never hard-code splits inline.
%------------------------------
% Weights: \Wl{l} = W^{(l)},  \Wsig{l}{\sigma} = W^{(l,\sigma)},  \posneg{W}{+} = [W]^+
\newcommand{\Wl}[1]{W^{(#1)}}
\newcommand{\Wsig}[2]{W^{(#1,#2)}}
\newcommand{\posneg}[2]{[#1]^{#2}}
% Activations: \actl{l} = a^{(l)},  \act{l}{\sigma} = a^{(l,\sigma)}
\newcommand{\actl}[1]{a^{(#1)}}
\newcommand{\act}[2]{a^{(#1,#2)}}
% Pre-activations: \preactl{l} = z^{(l)},  \preact{l}{\sigma} = z^{(l,\sigma)}
\newcommand{\preactl}[1]{z^{(#1)}}
\newcommand{\preact}[2]{z^{(#1,#2)}}
% Decomposition functions (optional argument = parameter)
\newcommand{\fmix}[1][\eta]{f_{\mathrm{mix}}^{\,#1}}
\newcommand{\fstop}{f_{\mathrm{stop}}}
\newcommand{\fsoft}[1][\theta]{f_{\mathrm{soft}}^{\,#1}}
\newcommand{\freg}[1][\Omega]{f_{\mathrm{reg}}^{\,#1}}
% Game theory: players P^+/P^-, generic p,q; opponent \opp{p} = p'
\newcommand{\opp}[1]{#1'}
\newcommand{\polstar}{\pi^{\star}}
\newcommand{\Vp}[2]{V_{#1}(#2)}   % game value V_p(s)
\newcommand{\Rp}[2]{R_{#1}(#2)}   % payoff R_p(s)

% Redefine mathcal to use mathscr globally
\renewcommand{\mathcal}{\mathscr}

\newenvironment{proofref}[2]{%
  \begin{proof}[Proof of #1~\ref{#2}]%
}{%
  \end{proof}%
}

\newcommand{\red}[1]{{\color{red}#1}}
%% \todo is already provided by the todonotes package; we define inline annotation variants
\newcommand{\todoinline}[1]{{\color{red}#1}}
\newcommand{\TODO}[1]{\textbf{\color{red}[TODO: #1]}}
%------------------------------
% COLORS
%------------------------------
\definecolor{dbta_blue}{RGB}{51,102,153}
\definecolor{tub_red}{RGB}{197,14,31}
%------------------------------
% XPARSE
%------------------------------
% O is an optional argument and {} is the default, m is the must-have argument
\NewDocumentCommand{\pder}{ O{} O{} m }{\dfrac{\partial^{#2}#1}{\partial#3^{#2}}}				% a short command for partial derivatives
\NewDocumentCommand{\roundbrack}{ m }{\left(#1 \right)} % a short command for round brackets
\NewDocumentCommand{\squarebrack}{ m }{\left[#1 \right]}% a short command for squared brackets
\NewDocumentCommand{\ind}{m O{}}{\mathrm{#1}\mathnormal{,#2}}			% a short command for indices
%------------------------------
% CODE
%------------------------------
\renewcommand{\lstlistingname}{Code}	% rename listings
%------------------------------
% ENUMERATE ENVIRONMENT
%------------------------------
\renewcommand{\labelenumi}{\arabic{enumi})} 				% 1)
\renewcommand{\labelenumii}{\arabic{enumi}.\arabic{enumii}}	% 1.1)
\renewcommand{\labelenumiii}{\arabic{enumi}.\arabic{enumii}.\arabic{enumiii}}	% 1.1.1)
\setlist[enumerate,1]{%
	itemsep=0pt,
	parsep=0pt,
	topsep=2pt,
	labelindent=\parindent,
}%
%------------------------------
% ITEMIZE ENVIRONMENT
%------------------------------
\setlist[itemize,1]{%
	label={--},
	itemsep=0pt,
	parsep=0pt,
	topsep=2pt,
	labelindent=\parindent,
}%
%------------------------------
% CAPTION REDEFINITIONS
%------------------------------
\addto\captionsenglish{
\renewcommand{\figurename}{Fig.}
\renewcommand{\tablename}{Tab.}
\renewcommand{\nomname}{List of Symbols}
\renewcommand{\lstlistlistingname}{List of Codes}
\def\indexname{Index}
}
%------------------------------
% ADD TOCLIST ALGORITHM
%------------------------------
\addtotoclist[algorithm2e]{loa}
\setuptoc{loa}{totoc}
\newcommand*{\listofloaname}{\listalgorithmcfname}
\renewcommand{\listofalgorithms}{\listoftoc{loa}}
%------------------------------
% AUTOREF ENGLISH
%------------------------------
\addto\extrasenglish{%
	\def\chapterautorefname{Chap.}%
	\def\sectionautorefname{Sec.}%
	\def\subsectionautorefname{Sec.}%
	\def\subsubsectionautorefname{Sec.}%
	\def\figureautorefname{Fig.}%
	\def\tableautorefname{Tab.}%
	\def\equationautorefname~#1\null{Eq.~(#1)\null}%
	\def\theoremenvautorefname{Theorem}%
	\def\lemmaenvautorefname{Lemma}%
	\def\remarkenvautorefname{Remark}%
	\def\definitionautorefname{Def.}%
	\def\algorithmautorefname{Alg.}%
	\def\assumptionautorefname{Ass.}%
	\def\chapterautorefnamelong{Chapter}%
	\def\sectionautorefnamelong{Section}%
	\def\subsectionautorefnamelong{Section}%
	\def\subsubsectionautorefnamelong{Section}%
	\def\figureautorefnamelong{Figure}%
	\def\tableautorefnamelong{Table}%
	\def\equationautorefnamelong~#1\null{Equation~(#1)\null}%
	\def\theoremenvautorefnamelong{Theorem}%
	\def\lemmaenvautorefnamelong{Lemma}%
	\def\remarkenvautorefnamelong{Remark}%
	\def\definitionautorefnamelong{Definition}%
	\def\algorithmautorefnamelong{Algorithm}%
}
%------------------------------
% AUTOREF (LONG FORM)
%------------------------------
\newcommand{\Autoref}[1]{%
  \begingroup%
  \renewcommand{\chapterautorefname}{\chapterautorefnamelong}%
  \renewcommand{\sectionautorefname}{\subsectionautorefnamelong}%
  \renewcommand{\subsectionautorefname}{\subsectionautorefnamelong}%
  \renewcommand{\subsubsectionautorefname}{\subsubsectionautorefnamelong}%
  \renewcommand{\figureautorefname}{\figureautorefnamelong}%
  \renewcommand{\tableautorefname}{\tableautorefnamelong}%
  \renewcommand{\equationautorefname}{\equationautorefnamelong}%
  \renewcommand{\theoremenvautorefname}{\theoremenvautorefnamelong}%
  \renewcommand{\lemmaenvautorefname}{\lemmaenvautorefnamelong}%
  \renewcommand{\remarkenvautorefname}{\remarkenvautorefnamelong}%
  \renewcommand{\definitionautorefname}{\definitionautorefnamelong}%
  \renewcommand{\algorithmautorefname}{\algorithmautorefnamelong}%
  \renewcommand{\assumptionautorefname}{\assumptionautorefnamelong}%
  \autoref{#1}%
  \endgroup%
}
%------------------------------
% REFERENCES AND CITATIONS
%------------------------------
\renewcaptionname{english}{\bibname}{List of References}
\renewcommand*{\bibfont}{\small}	% smaller font for the bibliography (saves pages)
\setlength\bibitemsep{0pt}		% no extra vertical gap between references
\setlength{\bibhang}{3ex}
\DeclareFieldFormat[			% no quotation marks in citetitle
	article,
	inbook,
	incollection,
	inproceedings,
	patent,
	thesis,
	unpublished
]{citetitle}{#1\midsentence\isdot}
\DeclareFieldFormat[			% no quotation marks in title
	article,
	inbook,
	incollection,
	inproceedings,
	patent,
	thesis,
	unpublished
]{title}{#1\midsentence\isdot}
\DeclareFieldFormat[			% remove capitalization in title and use grammar of a sentence instead
article,
]{titlecase:title}{\MakeSentenceCase*{#1}}
\DeclareFieldFormat[			% no quotation marks in title
	inproceedings,
]{booktitle}{\itshape #1}
\renewcommand*{\nameyeardelim}{\addcomma\space}			% comma between author and year
\DeclareFieldFormat{pages}{#1}	% no S. or pp. before pages
\DeclareFieldFormat{type}{\itshape #1}	% print bachelor or master thesis in italics
\DefineBibliographyStrings{english}{% for backreferencing in list of references
  backrefpage = {cit. on p.},
  backrefpages = {cit. on pp.},
}
\renewbibmacro{in:}{%			% suppress in: before name of journal, book, etc.
  \ifboolexpr{%
     test {\ifentrytype{article}}%
     or
     test {\ifentrytype{book}}%
     or
     test {\ifentrytype{inproceedings}}%
  }{}{\printtext{\bibstring{in}\intitlepunct}}%
}
\DefineBibliographyStrings{english}{% before URL english
urlseen = {last access},
}
\renewbibmacro*{volume+number+eid}{	% how volume and issue are printed in list of references
  \printfield{volume}
  \printfield[parens]{number}%
  \setunit{\addcomma\space}%
  \printfield{eid}
}
\DeclareDelimFormat[bib]{nametitledelim}{\addcolon\space} % use colon after year in list of references
\DeclareDelimFormat{finalnamedelim}{\addcomma\space}	  % delimiter between authors
\DeclareDelimFormat[textcite]{finalnamedelim}{			  % delimiter for textcite
  \addspace\&\space}
\DeclareDelimFormat[parencite]{finalnamedelim}{			  % delimiter for parencite
  \addspace\bibstring{and}\space}

\renewbibmacro*{begentry}{%
	\iffieldundef{shorthand}{}{%
		\printfield{shorthand}\addspace ---\addspace
	}
}
\DeclareSortingTemplate{nyt}{%		sorting template for shorthand
	\sort{%
		\field{presort}%
	}%
	\sort[final]{%
		\field{sortkey}%
	}%
	\sort{%
		\field{shorthand}%
		\field{author}%
		\field{editor}%
		\field{title}%
	}%
	\sort{%
		\field{year}
	}%
	\sort{%
		\field{title}%
	}%
}

\renewbibmacro*{textcite}{%			redefines textcite so that shorthands also work with it
    \ifnameundef{labelname}{%
    	\iffieldundef{shorthand}{%
    		\usebibmacro{cite:label}%
            \setunit{%
                \global\booltrue{cbx:parens}%
                \printdelim{nonameyeardelim}\bibopenparen%
            }%
            \ifnumequal{%
            	\value{citecount}%
            }{1}{%
            	\usebibmacro{prenote}%
            }{}%
            \usebibmacro{cite:labeldate+extradate}%
        }{%
        	\usebibmacro{cite:shorthand}%
        }}{%
        \iffieldundef{shorthand}{%
        	\printnames{labelname}%
            \setunit{%
                \global\booltrue{cbx:parens}%
                \printdelim{nameyeardelim}\bibopenparen%
            }%
           	\ifnumequal{%
           		\value{citecount}%
           	}{1}{%
           		\usebibmacro{prenote}%
           	}{}%
            \usebibmacro{citeyear}%
        }{%
        	\ifciteseen{%
        		\usebibmacro{cite:shorthand}%
        	}{%
        		\printnames{labelname}%
                \setunit{%
                    \global\booltrue{cbx:parens}%
                    \printdelim{nameyeardelim}\bibopenparen%
                }%
                \ifnumequal{%
                	\value{citecount}%
                }{1}{%
                	\usebibmacro{prenote}%
               	}{}%
                \iffieldundef{labelyear}{%
                	\usebibmacro{cite:label}%
                    \setunit{%
                    	\addcomma%
                    	\addspace%
                    }%
                    \iffieldundef{shorthandintro}{%
	                	\printtext{\bibstring{citedas}}%
	                }{%
	                	\printtext{\printfield{shorthandintro}}%
	                }%
	                \addspace%
                    \usebibmacro{cite:shorthand}%
                }{%
                	\usebibmacro{cite:labeldate+extradate}%
                    \setunit{%
                    	\addcomma%
                    	\addspace%
                    }%
                    \iffieldundef{shorthandintro}{%
	                    \printtext{\bibstring{citedas}}%
	                }{%
	                	\printtext{\printfield{shorthandintro}}%
	                }%
	                \addspace%
                    \usebibmacro{cite:shorthand}}}}}}%
%------------------------------
% ONLY FOR DISSERTATIONS
%------------------------------
\DeclareBibliographyCategory{inbib}	% a specific category for all citations in the text, NOT in the publication list
\newcommand{\printpublication}[1]{\AtNextCite{\defcounter{maxnames}{99}}\fullcite{#1}}						% print all authors in case the \printpublication command is used
\makeatletter
\AtEveryCitekey{%
  \ifcsstring{blx@delimcontext}{fullcite}	% if a citation is made with the fullcite command, do not show it in the references
    {}
    {\addtocategory{inbib}{\thefield{entrykey}}}}	% if it is normally cited later, add it to the references
\makeatother
%------------------------------
% INTERNET
%------------------------------
\renewcommand*{\UrlFont}{\upshape}
%------------------------------
% HYPERREF
%------------------------------
\hypersetup{%
   	bookmarksnumbered=true,		% bookmarks are numbered in PDF
   	bookmarksopen=false,			% bookmarks are closed when PDF is opened
   	colorlinks=true,			% colored text instead of boxed links (no ugly borders)
   	linkcolor={blue!45!black},	% internal links (sections, refs, equations)
   	citecolor={green!35!black},	% citations
   	urlcolor={blue!55!black},	% URLs
   	hypertexnames=true, 			% use guessable names for links
   	linktocpage, 				% use page numbers as links in table of contents
   	pdfpagelayout=SinglePage, 	% automatically open PDF as single page
}%
%------------------------------
% MAKEINDEX
%------------------------------
\makenomenclature				% generate an .nlo file for the List of Symbols
\providecommand{\pfill}{\dotfill}	% page-number fill used by myindexstyle.ist (absent from newer nomencl)
\newcommand{\indextitle}{Index}
\makeindex[
	columns=2,					% number of columns
	intoc=true,					% add to List of Contents
	options= -s bibliography/myindexstyle.ist,	% index style file
	title={\indextitle},	    		% title of the Index
	]							% generate an .idx file for the Index
%------------------------------
% VERTICAL SPACING / LAYOUT (fewer pages, no stretched-out gaps)
%------------------------------
% Do not stretch the page to fill it: prevents pages that are mostly one long
% vertical gap (KOMA/twoside otherwise uses \flushbottom).
\raggedbottom
% Let floats occupy more of a page before spilling, and allow more floats.
\renewcommand{\floatpagefraction}{0.75}
\renewcommand{\textfraction}{0.08}
\setcounter{totalnumber}{6}
% Tighten space around headings (KOMA). Negative beforeskip keeps the following
% paragraph unindented; the magnitude is the space above the heading.
\RedeclareSectionCommand[beforeskip=-2.0ex plus -1ex minus -.2ex, afterskip=1.0ex plus .2ex]{section}
\RedeclareSectionCommand[beforeskip=-1.7ex plus -.7ex minus -.2ex, afterskip=0.7ex plus .15ex]{subsection}
\RedeclareSectionCommand[beforeskip=-1.4ex plus -.5ex minus -.2ex, afterskip=0.5ex plus .1ex]{subsubsection}
% Trim the large blank band above KOMA chapter titles.
\RedeclareSectionCommand[beforeskip=0pt, afterskip=1.5ex]{chapter}
% This template writes \paragraph entries to the ToC regardless of tocdepth; drop them
% (and \subparagraph) so the contents list stays at chapter/section/subsection.
\let\dbtaOrigAddcontentsline\addcontentsline
\renewcommand{\addcontentsline}[3]{%
	\ifstrequal{#2}{paragraph}{}{\ifstrequal{#2}{subparagraph}{}{%
		\dbtaOrigAddcontentsline{#1}{#2}{#3}}}%
}
% Slightly tighten space around display math and theorem environments.
\setlength{\abovedisplayskip}{4pt plus 2pt minus 2pt}
\setlength{\belowdisplayskip}{4pt plus 2pt minus 2pt}
\setlength{\abovedisplayshortskip}{2pt plus 2pt}
\setlength{\belowdisplayshortskip}{3pt plus 2pt minus 2pt}

%------------------------------
% OWN COMMANDS
%------------------------------
% You can use this area for additional own commands
